\documentclass[11pt]{article}
\usepackage[margin=1in]{geometry}
\usepackage{amsmath, amssymb, amsthm}
\usepackage{algorithm, algpseudocode}
\usepackage{graphicx}
\usepackage{hyperref}

\title{Project Euler Problem 458:\\Permutations of Project}
\author{Solution}
\date{}

\newtheorem{theorem}{Theorem}
\newtheorem{lemma}{Lemma}

\begin{document}
\maketitle

\section{Problem Statement}

Consider the word \texttt{"project"} with 7 distinct letters $\mathcal{A} = \{p, r, o, j, e, c, t\}$.

Let $T(n)$ be the number of strings of length $n$ over $\mathcal{A}$ that do \emph{not} contain any
permutation of \texttt{"project"} as a contiguous substring of length 7 (i.e., no window of 7
consecutive characters contains all 7 distinct letters).

Given: $T(7) = 7^7 - 7! = 823543 - 5040 = 818503$.

\textbf{Goal:} Find $T(10^{12}) \bmod 10^9$.

\section{State Machine Model}

\subsection{State Definition}

Define state $s \in \{0, 1, 2, \ldots, 7\}$ as the length of the longest suffix of the string
built so far consisting of mutually distinct characters.

\begin{itemize}
    \item State 0: empty string or the last character creates no distinct suffix (initial state).
    \item State $s$ for $1 \le s \le 6$: the last $s$ characters are all distinct,
          but any extension to $s+1$ consecutive distinct chars has not yet occurred at this tail.
    \item State 7: \textbf{absorbing} --- a permutation of ``project'' has appeared.
\end{itemize}

\subsection{Transitions}

From state $s$ ($0 \le s \le 6$), adding a character $c \in \mathcal{A}$:

\begin{enumerate}
    \item \textbf{$c$ is new} (not among the last $s$ distinct chars): There are $7 - s$ such
          characters. Transition: $s \to s + 1$.
    \item \textbf{$c$ duplicates the $i$-th character from the right} ($1 \le i \le s$): There is
          exactly 1 such character for each $i$. The new longest suffix of distinct characters
          has length $i$ (the last $i-1$ characters before $c$ are still distinct from each other
          and from $c$, plus $c$ itself). Transition: $s \to i$.
\end{enumerate}

Total transitions from state $s$: $(7 - s) + s = 7$. \checkmark

From state 7: all 7 characters lead back to state 7 (absorbing).

\subsection{Transition Matrix}

The $8 \times 8$ transition matrix $\mathbf{M}$ with $M_{ij}$ = number of characters causing
transition from state $j$ to state $i$:

\[
\mathbf{M} = \begin{pmatrix}
0 & 0 & 0 & 0 & 0 & 0 & 0 & 0 \\
7 & 1 & 1 & 1 & 1 & 1 & 1 & 0 \\
0 & 6 & 1 & 1 & 1 & 1 & 1 & 0 \\
0 & 0 & 5 & 1 & 1 & 1 & 1 & 0 \\
0 & 0 & 0 & 4 & 1 & 1 & 1 & 0 \\
0 & 0 & 0 & 0 & 3 & 1 & 1 & 0 \\
0 & 0 & 0 & 0 & 0 & 2 & 1 & 0 \\
0 & 0 & 0 & 0 & 0 & 0 & 1 & 7
\end{pmatrix}
\]

\section{Matrix Exponentiation}

Starting from state vector $\mathbf{v}_0 = (1, 0, 0, 0, 0, 0, 0, 0)^\top$ (state 0):
\[
\mathbf{v}_n = \mathbf{M}^n \, \mathbf{v}_0
\]

The answer is:
\[
T(n) = \sum_{s=0}^{6} (\mathbf{M}^n)_{s,0} \pmod{10^9}
\]

Matrix exponentiation computes $\mathbf{M}^n$ in $O(8^3 \log n)$ operations, which for
$n = 10^{12}$ requires only about $8^3 \times 40 \approx 20{,}480$ multiplications modulo $10^9$.

\section{Verification}

For $n = 7$:
\begin{itemize}
    \item Total strings: $7^7 = 823{,}543$.
    \item Strings containing a permutation: exactly the $7! = 5{,}040$ permutations themselves
          (since the string length equals the alphabet size).
    \item $T(7) = 823{,}543 - 5{,}040 = 818{,}503$. \checkmark
\end{itemize}

\section{Complexity}

\begin{itemize}
    \item Matrix size: $8 \times 8$ (constant).
    \item Exponentiation: $O(\log n)$ matrix multiplications.
    \item Each multiplication: $O(8^3) = O(512)$ modular multiplications.
    \item \textbf{Total:} $O(512 \cdot \log_2(10^{12})) \approx O(512 \cdot 40) = O(20{,}480)$.
\end{itemize}

This runs in microseconds.

\section{Visualization}

\begin{figure}[h]
\centering
\IfFileExists{visualization.png}{\includegraphics[width=0.95\textwidth]{visualization.png}}{\fbox{Visualization unavailable in this archive}}
\caption{Top-left: Transition matrix heatmap. Top-right: Fraction $T(n)/7^n$ of valid strings.
Bottom-left: State probability distributions. Bottom-right: Growth comparison of $T(n)$ vs $7^n$.}
\end{figure}

\section{Answer}

\[
\boxed{423341841}
\]

\end{document}
